\section{General definitions}
\label{sec:am-vector-potential-split}

Following Ref.~\cite{bliokh_conservation_2014}, we introduce the standard
separation of electromagnetic angular momentum into orbital (OAM) and
spin (SAM) contributions through the vector potential.  We use SI units
and consider a source-free radiation field, so that
$\boldsymbol{\nabla}\cdot{\bf E}=0$.  All quantities in this section are
instantaneous fields evaluated at $(t,{\bf r})$, with ${\bf r}={\bf x}$
measured from the chosen origin.

\subsection{Total angular momentum and the vector potential}

The kinetic momentum density and its angular-momentum density are
\begin{align}
  {\bf g}&=\epsilon_0{\bf E}\times{\bf B},\qquad
  {\bf j}^{\mathrm{kin}}={\bf r}\times{\bf g},\\
  {\bf J}_V&=\int_V d^3r\,{\bf j}^{\mathrm{kin}}.
  \label{eq:am-general-total-integral}
\end{align}
We choose the transverse vector potential in the radiation gauge,
\begin{align}
  \boldsymbol{\nabla}\cdot{\bf A}=0,\qquad
  \Phi=0,\qquad
  {\bf E}=-\partial_t{\bf A},\qquad
  {\bf B}=\boldsymbol{\nabla}\times{\bf A}.
  \label{eq:am-general-radiation-gauge}
\end{align}
The transverse potential is fixed by the fields and suitable boundary
conditions.  This gauge choice specifies the orbital and spin densities
used below; expressions involving an arbitrary, unfixed potential would
in general depend on the gauge.

Substituting ${\bf B}=\boldsymbol{\nabla}\times{\bf A}$ gives, in components,
\begin{align}
  g_k=\epsilon_0\sum_m E_m\partial_k A_m
  -\epsilon_0\sum_m E_m\partial_m A_k.
  \label{eq:am-general-momentum-potential}
\end{align}
Using $\partial_m E_m=0$ and the product rule,
\begin{align}
  \varepsilon_{ijk}r_j E_m\partial_m A_k
  =\partial_m\left[E_m\varepsilon_{ijk}r_j A_k\right]
  -\varepsilon_{imk}E_m A_k,
\end{align}
where repeated spatial indices are summed.  Hence the local decomposition is
\begin{highlightbox}
  \begin{align}
    j_i^{\mathrm{kin}}
    =\mathcal L_i+\mathcal S_i
    -\epsilon_0\partial_m\left[E_m({\bf r}\times{\bf A})_i\right].
    \label{eq:am-general-local-split}
  \end{align}
\end{highlightbox}

\subsection{Orbital and spin contributions}

The canonical orbital and spin angular-momentum densities are
\begin{importantbox}
  \begin{align}
    \boldsymbol{\mathcal L}
    &=\epsilon_0\sum_{m=1}^{3}E_m
      ({\bf r}\times\boldsymbol{\nabla})A_m,
    \label{eq:am-general-orbital-density}\\
    \boldsymbol{\mathcal S}
    &=\epsilon_0{\bf E}\times{\bf A}.
    \label{eq:am-general-spin-density}
  \end{align}
\end{importantbox}
The orbital term contains the spatial rotation operator acting on each
Cartesian component of the potential and depends on the chosen origin.
The spin term contains no explicit position vector and describes the
polarization contribution.  Here $\boldsymbol{\mathcal S}$ denotes spin
angular-momentum density, distinct from the Poynting vector ${\bf S}$.

Integrating Eq.~\eqref{eq:am-general-local-split} over a volume $V$ gives
\begin{align}
  {\bf J}_V
  ={\bf L}_V+{\bf S}_V^{\mathrm{spin}}
  -\epsilon_0\oint_{\partial V}dA\,
    ({\bf E}\cdot{\bf n})({\bf r}\times{\bf A}),
  \label{eq:am-general-split-boundary}
\end{align}
where ${\bf n}$ is the outward unit normal,
${\bf L}_V=\int_V d^3r\,\boldsymbol{\mathcal L}$, and
${\bf S}_V^{\mathrm{spin}}=\int_V d^3r\,\boldsymbol{\mathcal S}$.
For fields localized sufficiently rapidly at infinity, or for boundary
conditions that make this surface integral vanish, the total angular
momentum therefore separates as
\begin{importantbox}
  \begin{align}
    {\bf J}={\bf L}+{\bf S}^{\mathrm{spin}},\qquad
    {\bf L}=\epsilon_0\int d^3r\sum_{m=1}^{3}
      E_m({\bf r}\times\boldsymbol{\nabla})A_m,\\
    {\bf S}^{\mathrm{spin}}
    =\epsilon_0\int d^3r\,{\bf E}\times{\bf A}.
    \label{eq:am-general-integrated-spin-orbital-split}
  \end{align}
\end{importantbox}
Locally, the canonical sum
$\boldsymbol{\mathcal L}+\boldsymbol{\mathcal S}$ and the kinetic density
${\bf r}\times{\bf g}$ differ by the divergence in
Eq.~\eqref{eq:am-general-local-split}.  For a finite integration region,
the corresponding surface term must be retained unless it vanishes.

\subsection{Components along the detector axis}

For the axis $Oz\equiv Ox_3$, the densities used in the following sections
are
\begin{highlightbox}
  \begin{align}
    \mathcal L_z
    &=\epsilon_0\sum_{m=x,y,z}E_m
      (x\partial_y-y\partial_x)A_m,\\
    \mathcal S_z&=\epsilon_0(E_x A_y-E_y A_x).
    \label{eq:am-general-axial-spin-orbital-densities}
  \end{align}
\end{highlightbox}
These are the time-domain definitions underlying the spectral expressions.
With the Fourier convention used in this document,
$\tilde{\bf A}=\tilde{\bf E}/(i\omega)$ for $\omega>0$.
